\section{Symbolic and numerical reproducibility}
\label{app:reproducibility}

All reported symbolic identities and numerical diagnostics are generated by
version-controlled source files; none of the displayed coefficients is
entered only in the manuscript.

\subsection{Symbolic certificates}

The module
\begin{center}
\texttt{src/canard\_control/final\_two\_module.py}
\end{center}
constructs the equilibrium, modal bases, layer matrices, projected atomic
measures, and current-state Jacobian.  Its tests verify:
\begin{enumerate}[label=\textup{(\roman*)}]
\item biorthogonality and the projector identities;
\item the fold curvature and exact slow critical line;
\item strict layer positivity on the declared \(\eta\)-box;
\item invariance of the complete projected delay measure and the signed
transverse forcing;
\item the exact algebraic center chain of the current-state singular
matrix.
\end{enumerate}
The current-state chain is used only as algebraic orientation; the proof
does not infer an RFDE spectral gap from it.

The module
\begin{center}
\texttt{src/canard\_control/final\_model\_blowup.py}
\end{center}
starts from \cref{eq:physical-model}, substitutes
\cref{eq:blowup-scaling}, projects with \(\ell^T,m^T\), and reconstructs
the original vector equation.  It verifies exact zero residuals for all four
chart equations and for the two reconstructed vector equations.  Finally,
\begin{center}
\texttt{src/canard\_control/nonlocal\_graph\_jet.py}
\end{center}
performs polynomial division by \(\del\) after the stable shift, verifies a
zero remainder rather than assuming divisibility, and derives
\cref{eq:Q1,eq:G0,eq:H1-eta,eq:transverse-return-coefficient}.

\subsection{Numerical diagnostic}

The numerical implementation is
\begin{center}
\texttt{src/canard\_control/exact\_chart\_threshold\_diagnostic.py},
\end{center}
with command-line driver
\begin{center}
\texttt{experiments/exact\_chart\_threshold\_diagnostic.py}.
\end{center}
It integrates each method-of-steps segment with
\texttt{scipy.integrate.solve\_ivp(method="Radau")}.  The reported run
uses relative tolerance \(2\times10^{-9}\), absolute tolerance
\(2\times10^{-11}\), and maximum internal step \(0.08\).  A Brent solve
uses \texttt{xtol=rtol=}\(2\times10^{-10}\); its bracket begins at the
zero-redistribution root and doubles until the finite energy gap changes
sign.

At \((\del,S)=(0.0025,5)\), the three tolerance pairs
\[
 (2\times10^{-9},2\times10^{-11}),\quad
 (5\times10^{-10},5\times10^{-12}),\quad
 (10^{-10},10^{-12})
\]
produce central quotients
\[
 -0.2036173761433,\quad
 -0.2036173733339,\quad
 -0.2036173737807,
\]
a spread of \(2.81\times10^{-9}\).  With the first tolerance pair and
maximum steps \(0.08,0.04,0.02\), the corresponding spread is
\(3.16\times10^{-9}\).  The much larger discrepancy from
\cref{eq:numerical-target} is therefore asymptotic/history/section error,
not ordinary solver-tolerance variation.

The machine-readable archive
\begin{center}
\texttt{experiments/results/exact\_chart\_threshold\_convergence.json}
\end{center}
contains all full-precision roots, one-sided and central quotients,
residuals, tolerances, software versions, and source hashes.  The paper
figure is generated directly from that file by
\begin{center}
\texttt{manuscript/jns/figures/generate\_figures.py}.
\end{center}

\subsection{Reproduction commands}

From the repository root, the symbolic/unit tests are run with
\begin{verbatim}
python -m pytest -q
\end{verbatim}
and the paper assets and PDF with
\begin{verbatim}
make -C manuscript/jns figures
make -C manuscript/jns paper
\end{verbatim}
The manuscript build uses Tectonic; the numerical and paper extras install
NumPy, SciPy, SymPy, mpmath, and ReportLab.  The repository records the exact
command and environment used for the reported numerical table.  A versioned
code-and-data archive accompanies the manuscript as supplementary material;
a public archival copy with a persistent identifier will be released upon
acceptance.
